\section{Công trình liên quan}
\subsection{Các công trình khoa học liên quan}
Các nghiên cứu về hệ thống khuyến nghị du lịch (Tourism Recommender System – TRS) chủ yếu tập trung vào khai thác dữ liệu cộng đồng nhằm nâng cao tính cá nhân hóa và đảm bảo tính chân thực của thông tin.

Chẳng hạn, Khan et al. (2025) đã áp dụng mô hình BERT để phân tích ngữ nghĩa sâu từ 3.013 bài báo khoa học giai đoạn 2000–2024, kết hợp kỹ thuật giảm chiều UMAP và phân cụm HDBSCAN, từ đó xác định 16 tham số chính được nhóm thành bốn nhóm tham số (macro-parameters): Personalized Tourism (tập trung vào cá nhân hóa hành trình), Sustainability, Health \& Resource Awareness (bền vững và nhận thức về tài nguyên sức khỏe), Adaptability \& Crisis Management (thích ứng và quản lý khủng hoảng), Social Impact \& Cultural Heritage (tác động xã hội và di sản văn hóa). Phương pháp này khai thác dữ liệu do người dùng tạo (user-generated content) để giảm thiểu hiện tượng đánh giá giả (seeding) và nâng cao tính xác thực, đồng thời nhấn mạnh vai trò của AI trong phân tích dữ liệu lớn phục vụ khuyến nghị dựa trên đánh giá cộng đồng.

Tương tự, Patil và Patil (2024) đề xuất một kiến trúc TRS thông minh tích hợp AI, trong đó quá trình cá nhân hóa dựa trên sở thích và dữ liệu nhân khẩu học của người dùng. Kiến trúc gồm ba thành phần chính: user profiling (học từ lịch sử tương tác), content filtering (mô hình lai – hybrid), và trip planning (tối ưu hóa lịch trình). Phương pháp này sử dụng thuật toán K-Means để phân cụm dữ liệu và giải quyết vấn đề cold-start thông qua dữ liệu nhân khẩu học. Tuy nhiên, nghiên cứu chủ yếu tập trung vào nâng cao sự hài lòng của người dùng, chưa mở rộng sang hỗ trợ các nhà cung cấp dịch vụ (SMEs).

Ngoài ra, Caliciuri et al. (2024) đã thực hiện một khảo sát tổng quan, phân loại các nghiên cứu về khuyến nghị hành trình cá nhân hóa thành hai hướng chính: user satisfaction (bao gồm personalized và non-personalized dựa trên sở thích và ràng buộc của khách du lịch) và provider satisfaction (tối ưu hóa tài nguyên cung ứng). Khảo sát này nhấn mạnh vai trò của dữ liệu cộng đồng trong việc xây dựng lịch trình tùy chỉnh và đề xuất các hướng nghiên cứu tương lai như công bằng (fairness) và tối ưu hóa sử dụng tài nguyên.

Nhìn chung, các công trình hiện tại chủ yếu dừng ở mức lý thuyết và phân tích dữ liệu, chưa triển khai thành các hệ thống tích hợp đầy đủ trong thực tiễn. Đặc biệt tại Việt Nam, vẫn còn thiếu các hệ thống phù hợp với nhu cầu thế hệ Z (yêu cầu thông tin chân thực, cập nhật nhanh) cũng như công cụ để SMEs quảng bá dịch vụ và xác thực đánh giá theo thời gian thực. Khoảng trống nghiên cứu này là cơ hội quan trọng để phát triển giải pháp hỗ trợ du lịch số, nâng cao minh bạch và trải nghiệm người dùng.

\subsection{Các hệ thống/phần mềm hiện có}
\subsubsection{Wanderlog}
Wanderlog là một ứng dụng lập kế hoạch hành trình du lịch cá nhân, nổi bật với khả năng tích hợp bản đồ và chi phí ngay trong quá trình xây dựng lịch trình. Người dùng có thể thêm thủ công hoặc tìm kiếm các điểm đến, sau đó hệ thống tự động hiển thị trên bản đồ và cho phép tối ưu hóa tuyến đường. Wanderlog hỗ trợ chia sẻ lịch trình với bạn bè, đồng thời có chế độ đồng bộ hóa nhiều thiết bị. Ứng dụng cung cấp phiên bản miễn phí và bản Pro với nhiều tính năng nâng cao như xuất PDF, đồng bộ offline và hỗ trợ nhóm du lịch. Trang web chính thức tại \href{wanderlog.com}{wanderlog.com} cung cấp thông tin chi tiết về tính năng, hướng dẫn sử dụng và tải ứng dụng.

\subsubsection{Stippl}
Stippl là ứng dụng lên kế hoạch hành trình toàn diện với các tính năng quản lý ngân sách, checklist chuẩn bị và chia sẻ trải nghiệm du lịch. Người dùng có thể lên kế hoạch theo ngày, tính toán chi phí, quản lý danh sách công việc cần chuẩn bị và chia sẻ hành trình với cộng đồng Stippl. Điểm mạnh của Stippl là kết hợp yếu tố tài chính (budgeting) cùng tính năng xã hội, giúp khách hàng không chỉ đi du lịch mà còn chia sẻ và kết nối đam mê. Thông tin chi tiết có tại \href{stippl.io}{stippl.io}.

\subsubsection{GuideGeek}
GuideGeek là trợ lý du lịch AI do Matador Network phát triển. Điểm đặc biệt là hệ thống vận hành dưới dạng chatbot, có thể tích hợp trên các nền tảng phổ biến như Instagram, Messenger, WhatsApp. GuideGeek sử dụng AI để đưa ra gợi ý điểm đến, tip du lịch, thông tin ăn uống và hoạt động phù hợp với người dùng. Đây là một hướng đi khác biệt, tập trung vào sự tiện lợi và cá nhân hóa, thay vì lập lịch trình chi tiết. Website chính thức \href{guidegeek.com}{guidegeek.com} cung cấp bản demo chatbot và các case study ứng dụng thực tế.

\subsubsection{TripHunter AI Chatbot}
TripHunter AI Chatbot là sản phẩm tiên phong ứng dụng AI trong ngành du lịch Việt Nam. Chatbot này được tích hợp trực tiếp vào hệ thống TripHunter, hỗ trợ khách hàng tư vấn, cá nhân hóa gợi ý dịch vụ và trả lời tự động. Với khả năng xử lý ngôn ngữ tự nhiên (NLP), chatbot giúp nâng cao trải nghiệm khách hàng, tăng tính tương tác và giảm tải cho nhân viên hỗ trợ. Đây là bước đi quan trọng trong việc áp dụng công nghệ AI vào dịch vụ du lịch Việt Nam, đặc biệt là trong bối cảnh nhu cầu tư vấn và đặt dịch vụ trực tuyến ngày càng tăng.

\subsubsection{Visit Vietnam}
Visit Vietnam là "siêu ứng dụng" và nền tảng dữ liệu du lịch quốc gia, được phát triển bởi Tập đoàn Sun Group dưới sự bảo trợ của Chính phủ, Cục Du lịch Quốc gia Việt Nam (VNAT) và Hiệp hội Dữ liệu Quốc gia (NDA). Nền tảng này được định vị là hạ tầng số cốt lõi, phục vụ đồng thời ba đối tượng: cơ quan quản lý nhà nước, doanh nghiệp du lịch và du khách.

Visit Vietnam tích hợp công nghệ bản đồ số 3D/4D và thực tế ảo (AR/VR), cho phép du khách tương tác và trải nghiệm điểm đến từ xa. Điểm nổi bật nhất của hệ thống là khả năng vận hành dữ liệu theo thời gian thực (real-time data), giúp du khách theo dõi mức độ đông đúc tại điểm đến để tránh quá tải. Hệ thống tích hợp trợ lý ảo "AI Travel Assistant" hỗ trợ xây dựng lịch trình cá nhân hóa dựa trên thời gian, ngân sách và sở thích riêng biệt của từng người dùng \cite{vnexpress2025_dulich}. 

Bên cạnh đó, Visit Vietnam hợp tác với Visa để tích hợp phân tích hành vi chi tiêu và thanh toán số, tạo nên trải nghiệm "du lịch một chạm" liền mạch từ tìm kiếm đến đặt dịch vụ \cite{cafef2025}. Đối với doanh nghiệp, nền tảng cung cấp các công cụ quản trị thông minh (Dashboard) và báo cáo xu hướng thị trường chuyên sâu giúp tối ưu hóa vận hành. Dự kiến, nền tảng sẽ vận hành chính thức vào quý II/2026. Do đó, chúng tôi chỉ ghi nhận các tính năng công bố như một tham chiếu về xu hướng công nghệ (State-of-the-Art) mà chưa đưa vào bảng đánh giá so sánh chức năng thực tế.

\subsection{So sánh với hệ thống của nhóm}
Chúng tôi sẽ tiến hành so sánh hệ thống của mình với hệ thông khác trong ngành dựa trên các tiêu chí đánh giá và phân mức cấp độ chi tiết sau:
\begin{enumerate}
    \item Lập kế hoạch hành trình
        \begin{itemize}
            \item \textbf{Cấp 1}: Người dùng nhập thủ công điểm đến → hiển thị danh sách/bản đồ.
            \item \textbf{Cấp 2}: Tự động sắp xếp thứ tự điểm đến theo khoảng cách/thời gian.
            \item \textbf{Cấp 3}: Tối ưu đa tiêu chí (thời gian, chi phí, sở thích, phương tiện) bằng AI, cập nhật thời gian thực.
        \end{itemize}
    \item Cá nhân hóa \& Gợi ý thông minh
        \begin{itemize}
            \item \textbf{Cấp 1}: Lọc cơ bản (giá, loại hình dịch vụ).
            \item \textbf{Cấp 2}: Gợi ý dựa trên sở thích đã khai báo hoặc hành vi trước.
            \item \textbf{Cấp 3}: Recommendation Engine (ML/DL), NLP phân tích review, gợi ý động và cá nhân hóa sâu
        \end{itemize}
    \item Tích hợp dịch vụ du lịch (Booking \& Services)
        \begin{itemize}
            \item \textbf{Cấp 1}: Hiển thị thông tin dịch vụ (khách sạn, vé, tour).
            \item \textbf{Cấp 2}: Cho phép đặt chỗ/đặt vé trực tiếp trong app.
            \item \textbf{Cấp 3}: Thanh toán điện tử, đồng bộ đa dịch vụ (flight + hotel + activity) trên cùng một nền tảng.
        \end{itemize}
    \item Đánh giá, bình luận \& cộng đồng
        \begin{itemize}
            \item \textbf{Cấp 1}: Chỉ hiển thị review/bình luận từ nguồn sẵn có.
            \item \textbf{Cấp 2}: Người dùng được viết đánh giá, bình luận và nhận phản hồi.
            \item \textbf{Cấp 3}: Diễn đàn/Community, xác thực chống review ảo, gợi ý dựa trên cộng đồng.
        \end{itemize}
    \item Quản lý chi phí \& Ngân sách
        \begin{itemize}
            \item \textbf{Cấp 1}: Hiển thị chi phí cơ bản (giá dịch vụ).
            \item \textbf{Cấp 2}: Người dùng có thể nhập/lập ngân sách thủ công.
            \item \textbf{Cấp 3}: Tự động dự toán chi phí, gợi ý hành trình tối ưu theo ngân sách.
        \end{itemize}
    \item Trải nghiệm người dùng (UX/UI)
        \begin{itemize}
            \item \textbf{Cấp 1}: Giao diện cơ bản, chưa tối ưu di động.
            \item \textbf{Cấp 2}: Giao diện tối ưu cho di động, có bản đồ tương tác.
            \item \textbf{Cấp 3}: UI/UX thân thiện, hỗ trợ tùy chỉnh cao, bản đồ real-time, offline mode.
        \end{itemize}
    \item Hỗ trợ doanh nghiệp/đối tác
        \begin{itemize}
            \item \textbf{Cấp 1}: Chỉ hiển thị danh sách dịch vụ từ đối tác
            \item \textbf{Cấp 2}: Công cụ quản lý đơn hàng/tour cho doanh nghiệp.
            \item \textbf{Cấp 3}: Dashboard phân tích dữ liệu.
        \end{itemize}
\end{enumerate}

% --- BẢNG ĐÃ ĐƯỢC XOAY NGANG ĐỂ ĐẢM BẢO HIỂN THỊ ĐỦ ---
\clearpage
\begin{landscape}
\begin{table}[h!]
\centering
\footnotesize % Giảm cỡ chữ để có thêm không gian
\renewcommand{\arraystretch}{1.3} % Tăng khoảng cách dòng cho dễ đọc
\setlength{\tabcolsep}{4pt} % Giảm khoảng cách giữa các cột
\begin{tabularx}{\linewidth}{|>{\raggedright\arraybackslash}X|>{\centering\arraybackslash}X|>{\centering\arraybackslash}X|>{\centering\arraybackslash}X|>{\centering\arraybackslash}X|}
\hline
\textbf{Tiêu chí} & 
\textbf{Wanderlog} & 
\textbf{Stippl} & 
\textbf{GuideGeek} & 
% \textbf{Visit Vietnam} & 
% \textbf{Travelviet} & 
% \textbf{My Travel Map / Gody} & 
\textbf{TripHunter AI Chatbot} \\ 
\hline
Lập kế hoạch hành trình & 3 & 2 & 1 & 1 \\ 
\hline
Cá nhân hoá \& Gợi ý thông minh & 2 & 2 & 3 & 3 \\ 
\hline
Tích hợp dịch vụ du lịch & 2 & 1 & 1 & 2 \\ 
\hline
Đánh giá, bình luận \& cộng đồng & 2 & 2 & 1 & 2 \\ 
\hline
Quản lý chi phí \& Ngân sách & 3 & 3 & 1 & 1 \\ 
\hline
Trải nghiệm người dùng (UX/UI) & 3 & 2 & 2 & 2 \\ 
\hline
Hỗ trợ doanh nghiệp/đối tác & 1 & 1 & 1 & 2 \\ 
\hline
\end{tabularx}
\caption{Đánh giá các hệ thống theo tiêu chí chức năng (cấp độ 1–3).}
\label{tab:danh_gia_he_thong_landscape}
\end{table}
\end{landscape}
